\section{Discussion}
\label{sec:discussion}

\para{Vibe is a real and distinct signal.}
Our results provide strong evidence that LLM-generated vibe descriptions capture information genuinely different from what raw-text embeddings encode.
The low \ari (0.10--0.14), large effect size on neighbor content similarity (Cohen's $d{=}1.44$), and consistent results across all values of $k$ all point to the same conclusion: vibe space is not a noisy version of content space but a qualitatively different organization of web content.

\para{Why does the pipeline work?}
We hypothesize that the ``it's giving'' prompt activates a specific mode of cultural commentary in the LLM---one trained on internet discourse where aesthetic judgments are common.
The LLM compresses each page's content, layout cues (reflected in text structure), and cultural markers into a concise aesthetic summary.
The sentence transformer then embeds these summaries in a space where aesthetic similarity translates to geometric proximity.

\para{Implications for web navigation.}
The existence of a coherent vibe space suggests that ``browse by vibe'' is a feasible navigation modality.
A user could select a page they enjoy and ask for pages that ``give the same energy,'' surfacing content from entirely different domains.
This is complementary to topic-based search: where search answers ``what is this about?,'' vibe navigation answers ``what does this \emph{feel} like?''

\para{Limitations.}
Several limitations qualify our findings.
First, our sample of 497 pages, while sufficient for statistical testing, is far from web scale.
Whether vibe clusters remain coherent at millions of pages is an open question.
Second, we use pre-extracted text from \cfour rather than live HTML; visual design elements (colors, fonts, layout) contribute to vibe but are not captured by our text-only pipeline.
Third, we test only one LLM (\gptfour) with one prompt; different models or prompts may produce different vibe characterizations.
Fourth, we lack a user study: we demonstrate structural properties of vibe space but do not test whether humans find it useful for browsing.
Fifth, the \cfour corpus is predominantly English, and vibe characterizations may be culturally specific to English-language internet culture.
Finally, cost is a barrier to scale: at approximately \$0.01 per page, processing one billion pages would cost roughly \$10 million without optimization.

\para{Broader impact.}
A vibe-based navigation system could help users discover content outside their usual information bubbles---not by recommending opposing viewpoints (as diversity interventions typically do), but by connecting content that shares aesthetic character across domain boundaries.
However, vibes are subjective and culturally situated.
A system that foregrounds ``vibes'' risks encoding the aesthetic preferences of the LLM's training data, which skews toward English-language, Western internet culture.
Deploying such a system responsibly would require careful attention to whose vibes are represented and whose are marginalized.
